% 结构版本：D038 B 方案大纲（③终选）。4.1–4.2 定稿 V012（2026-09-26，W-C 窗口）；4.3–4.6 为骨架，小节标题已按 D038 同步，正文待各窗口起草（不随本轮生成）。
\section{场景 A 下的多核切图与调度模型}\label{sec:q1}

\subsection{场景 A 的建模描述}\label{sec:q1-desc}
% 状态：V012 定稿（2026-09-26 用户对话审定，两轮 GPT 评审逐条仲裁后写入，AI-13/A049）。素材：题面 1.5 场景 A、问题 1、附录 B.5/D.2；官方评估器 evaluation_validation.py（validate\_task\_order 合并查环）与 \_build\_scene\_a\_tasks 语义核验；无实验数据、不挂文献。要点：融合式开篇三段（建模对象／场景 A 执行语义／职责分界＋一句方法路线）；不重述 I/O 与交付（1.5 节与 TBL-001 已覆盖）；激活常量 100/1000 cycles 为 1.2.3 节延后项落点；FastEval 定位＝加速评估实现（措辞与大纲同步，D038 行“搜索加速代理”已改口径）。
问题一的核心权衡已在 2.2 节分析。本章在场景 A 下建立多核切图与调度模型：模型的直接对象即 1.5.1 节所述的三层决策——切图、核心分配与同核执行顺序，输出为切图方案与子图调度方案两个字段（1.2.3 节）；本章的工作是将这一决策问题形式化，并在官方评估语义下求解。

场景 A 无核间同步机制：每个子图独立构成一个 Task，子图与 Task 一一对应；Task 之间不保留核内缓存状态，跨子图数据一律经 DDR 中转、由评估程序自动插入的 COPY\_OUT 与 COPY\_IN 传递，不因两个子图是否位于同一核心而改变（1.2.3 节）。Task 整体激活：其最早激活时刻同时受同核顺序与依赖前驱两类约束——同核前一 Task 完成后至少等待 100 cycles，每个跨核前驱 Task 完成后至少等待 1000 cycles，实际激活时刻取全部约束的最大值。

参赛算法与官方评估流程的职责由此分界：算法直接控制的只有上述两个方案字段，Task 内部操作顺序、边界搬运的插入、缓存换入换出与多核并行执行均由官方评估程序按固定流程完成（1.2.3、1.4 节），方案优劣以其模拟输出为准；搜索过程中使用的任何中间评价量——候选构造与指派中的启发式代价、加速评估实现——均不改变这一口径（4.2.4 节）。在此范围内，本章按 2.5 节的两段式路线组织求解：4.2 节给出形式化模型与评价边界，4.3～4.5 节建立构造、指派与在线求解方法，4.6 节报告实验结果。

\subsection{多核切图与调度优化模型}\label{sec:q1-model}
% 状态：V012 定稿（同上）。要点：决策 D=(x,{π_k}) 与方案两字段一一对应（核心分配 φ 由序列归属导出；sgid 非连续、以标签集合 S=x(O_c) 为下标集）；可行域两级 𝒟plan/𝒟feas（与 1.3.1 两层框架同构），优化在 𝒟feas 上定义；五条静态约束，其中（5）增强图 Q_D=E_dep∪E_ord 无环为独立检查层（官方 validate\_task_order 合并查环，反例经代码核证）；三地板＝方案机制下界（L_bw/L_relay 方案相关，不作最优值估计、不参与搜索）；FastEval＝与官方流程逐段对应的加速评估实现（一致性证据留 4.6.1）。Needs-integration：符号清单入 06/3.2 TBL-003；附录 D 三地板推导挂接（theorems.tex 引理 8–10）待整合窗；大纲 4.2.4 行措辞已同轮同步。
本节将 4.1 节的建模范围落实为形式化优化模型：4.2.1 节给出决策表示，4.2.2 与 4.2.3 节分别给出目标函数（含可行方案的分级定义）与约束条件，4.2.4 节说明模型依据、评价边界与方案机制下界。

\subsubsection{决策变量}\label{sec:q1-vars}
输入计算图的记号沿用 1.2.2 节：$G=(V,E)$，张量节点集合 $T$、操作节点集合 $O$，张量大小 $s_t$，核内计算操作执行周期 $c_o$；核心数量为 $N$。参与划分的对象为核内操作节点（1.2.3 节），记其集合为 $O_{\mathrm{c}}$，即 $O$ 中除去 COPY\_IN 与 COPY\_OUT 两类搬运操作后的全体操作节点。一个完整决策由两部分构成。

（1）\textbf{切图}：映射 $x:\ O_{\mathrm{c}}\rightarrow\mathbb{Z}_{\ge 0}$，将每个核内操作节点指派到一个非负整数子图标识 sgid。sgid 仅作标识、不要求连续取值，故记实际使用的标签集合 $S=x(O_{\mathrm{c}})$；对 $s\in S$，非空原像
\[
O_s=\{\,o\in O_{\mathrm{c}}:\ x(o)=s\,\}
\]
即子图 $s$ 的操作集合。映射 $x$ 即切图方案字段 \texttt{node\_to\_subgraph}。

（2）\textbf{子图调度}：对每个核心 $k\in\{1,\dots,N\}$，一个由 sgid 组成的有穷序列 $\pi_k=(\pi_{k,1},\dots,\pi_{k,r_k})$，按执行顺序列出该核心分得的全部子图，序列可为空。全部序列合起来即子图调度方案字段 \texttt{core\_schedules}。

核心分配不是独立变量：子图 $s$ 所属的核心由其出现在哪个序列中确定，记 $\phi(s)=k$ 当且仅当 $s\in\pi_k$——两个输出字段正以此方式同时给出分配与顺序（1.2.3 节）。因此，决策 $D=(x,\{\pi_k\}_{k=1}^{N})$ 与两个方案字段一一对应，完整表示切图、核心分配与同核执行顺序三层决策（2.2 节）；约束条件（4.2.3 节）进一步剔除不合法的取值组合。

\subsubsection{目标函数}
模型以缩短总体任务执行时间为主要目标。对决策 $D$，记 $\mathrm{mk}(D)$ 为官方评估程序在场景 A 语义、核心数量 $N$ 下对相应方案完成多核模拟后输出的总体任务执行时间（1.4 节）；$\mathrm{mk}$ 没有解析表达式，其值由 4.1 节所述评估流程的离散事件模拟给出。满足 4.2.3 节全部静态条件的方案构成方案级合法集合 $\mathcal{D}_{\mathrm{plan}}$，其中能被官方评估程序成功执行完毕的方案构成可行集合 $\mathcal{D}_{\mathrm{feas}}\subseteq\mathcal{D}_{\mathrm{plan}}$（分级依据见 4.2.3 节末）。问题一的优化问题在 $\mathcal{D}_{\mathrm{feas}}$ 上定义：
\begin{equation}\label{eq:q1-object}
\min_{D\in\mathcal{D}_{\mathrm{feas}}}\ \mathrm{mk}(D)
\end{equation}

总额外数据搬运量记为 $\mathrm{add}(D)$（1.4 节），作为次要指标随方案一并报告，用于结果分析与方案对照；本文不构造执行时间与搬运量的加权目标，方案接受与最终优劣的判据均为 $\mathrm{mk}$。与搬运相关的静态特征（如跨子图张量字节，1.3.2 节）进入候选构造与指派的启发式代价以引导搜索方向（4.4 节），不改变评价口径。评价的高频调用开销及相应的加速实现见 4.2.4 节。

\subsubsection{约束条件与合法性}\label{sec:validation}
合法方案须满足题面规定的全部方案条件（1.2.3 节）。按决策 $D=(x,\{\pi_k\})$ 表述为五条：（1）（2）界定两层决策各自的完整性，（3）～（5）为依赖合法性约束。

（1）\textbf{完整且唯一的划分归属}。切图须覆盖全部核内操作节点，且每个节点恰属一个子图。由于 $x$ 以 $O_{\mathrm{c}}$ 全体为定义域，映射性质自动保证各非空原像两两不交、并集覆盖：
\begin{equation}
\bigsqcup_{s\in S} O_s = O_{\mathrm{c}}
\end{equation}
其中 $\bigsqcup$ 表示不交并。

（2）\textbf{子图调度的完整性与唯一性}。每个子图必须且只能被调度一次：全部核心的执行序列合并后，恰为标签集合 $S$ 的一个排列；个别核心的序列可为空（空核）。此条同时保证核心归属 $\phi$ 良定义——每个 sgid 恰在一个序列中出现。

（3）\textbf{子图依赖图无环}。将操作间依赖按划分收缩，两端分属不同子图的依赖构成子图间依赖边，得到子图依赖图 $Q=(S,\,E_{\mathrm{dep}})$（构建方式见 1.3.1 节）。输入计算图为 DAG 并不保证任意划分下 $Q$ 无环——相互依赖的操作被交叉归入不同子图时，子图间即可形成环形依赖。$Q$ 须为有向无环图，方案才存在合法的执行顺序。

（4）\textbf{同核顺序与依赖的一致性}。对 $E_{\mathrm{dep}}$ 中每条依赖边 $s\rightarrow s'$，若两个子图位于同一核心 $k$，则执行顺序须与依赖方向一致。以 $\mathrm{pos}_k(s)$ 表示子图 $s$ 在序列 $\pi_k$ 中的位次，要求
\begin{equation}
\mathrm{pos}_k(s) < \mathrm{pos}_k(s'),\qquad \forall\, s\rightarrow s'\in E_{\mathrm{dep}},\ \phi(s)=\phi(s')=k
\end{equation}

（5）\textbf{整体调度可行性}。将各核执行序列的相邻顺序边
\[
E_{\mathrm{ord}}=\{\,(\pi_{k,r},\,\pi_{k,r+1})\mid 1\le k\le N,\ 1\le r<r_k\,\}
\]
加入子图依赖图，要求增强图 $Q_D=(S,\ E_{\mathrm{dep}}\cup E_{\mathrm{ord}})$ 为有向无环图。该条件不能由（3）（4）共同推出：即使 $Q$ 无环且无任何同核依赖边被逆置，顺序边与跨核依赖边组合仍可能成环——例如，若子图依赖图含 $B\rightarrow C\rightarrow A$，且同核顺序规定 $A$ 先于 $B$，则新增顺序边 $A\rightarrow B$ 后形成 $A\rightarrow B\rightarrow C\rightarrow A$ 的环。官方评估程序在构造 Task 时将子图依赖与各核串行任务顺序合并检查无环，任一环出现即终止评估。

以上五条均可静态检查，覆盖官方评估入口中与决策方案有关的全部静态合法性检查（1.3.1 节），据此构成 $\mathcal{D}_{\mathrm{plan}}$ 的完整判据；本文求解流程直接采用该判据，不另设判定口径。$\mathcal{D}_{\mathrm{plan}}$ 与 $\mathcal{D}_{\mathrm{feas}}$ 的差来自评估流程的执行阶段：核内调度在 L1/UB 容量不足时自动插入缓存换入换出，且在单个张量超过缓存容量或不存在合法换出对象时任务失败、评估终止（题面附录 C.2）；此类执行可行性不在静态约束中显式建模。

\subsubsection{模型推导与评价边界}\label{sec:q1-evalbound}
约束条件的依据是题面对方案表示的硬性规定（1.2.3 节），不引入附加假设；目标函数直接取题面指定的主要指标（1.4 节）。模型不解析展开 $\mathrm{mk}$ 的内部构成：计算操作的管道占用、DDR 搬运的带宽争用、Task 激活的依赖等待与缓存换入换出等成本，全部由官方评估程序在模拟中结算；静态侧仅使用 1.3.2、1.3.3 节的输入字段与静态推导特征支撑候选构造，不以张量大小之和构造容量约束，也不静态近似换入换出（1.3.3 节）。由此，模型在形式上是“\textbf{静态约束集合＋黑盒目标函数}”的优化问题，求解范式为官方评估语义下的方案搜索（2.5 节）。

$\mathrm{mk}$ 的单次求值需运行完整的离散事件模拟，高频调用构成搜索的主要开销。为此，本文采用与官方场景 A 评估流程逐段对应的加速评估实现 \textbf{FastEval}，用于候选筛选与精修中的反复比较；其与官方评估的一致性验证见 4.6.1 节，论文正式报告的数值均由官方评估程序复评给出。

黑盒目标下，单个方案的质量还需要不依赖求解过程的解析参照。本文为场景 A 构造任意合法方案的\textbf{必要耗时下界}。记 $W_M$、$W_V$ 为全部核内计算操作按执行管道 PIPE\_M、PIPE\_V 分别累计的分管道计算量（1.3.2 节），$W_M^{(s)}$、$W_V^{(s)}$ 为子图 $s$ 内的对应量，$\beta_0$ 为原始计算图固有的 DDR 搬运字节数，$B_{\mathrm{ddr}}$ 为 DDR 总带宽（60 bytes/cycle，1.2.1 节），则对任意 $D\in\mathcal{D}_{\mathrm{feas}}$：
\begin{equation}\label{eq:q1-lb}
\begin{aligned}
\mathrm{mk}(D)\ \ge\ L(D)&=\max\{\,L_{\mathrm{load}},\ L_{\mathrm{bw}}(D),\ L_{\mathrm{relay}}(D)\,\},\\
L_{\mathrm{load}}&=\frac{\max(W_M,\,W_V)}{N},\qquad
L_{\mathrm{bw}}(D)=\frac{\beta_0+\mathrm{add}(D)}{B_{\mathrm{ddr}}},\\
L_{\mathrm{relay}}(D)&=\max_{s_1\rightarrow\cdots\rightarrow s_r\,\in\,Q}\Bigg(\sum_{l=1}^{r}\max\big(W_M^{(s_l)},\,W_V^{(s_l)}\big)+\sum_{l=1}^{r-1}\omega\big(s_l,\,s_{l+1}\big)\Bigg)
\end{aligned}
\end{equation}
三条下界各对应执行机制中的一类瓶颈：$L_{\mathrm{load}}$ 源于每核每管道严格串行，至少一个核心的管道累计负载被垫至平均水平之上；$L_{\mathrm{bw}}$ 源于全部 DDR 搬运共享总带宽——任一时刻的聚合服务速率不超过 $B_{\mathrm{ddr}}$，存在多条并发搬运时带宽按公平共享方式分配，而全部搬运须在总完成时间内完成；$L_{\mathrm{relay}}$ 源于依赖链上的逐段等待——链上各 Task 依序激活，$\omega(s,s')$ 按同核 100 cycles、跨核 1000 cycles 取值（4.1 节），子图内管道串行使 $\max(W_M^{(s)},W_V^{(s)})$ 构成其执行时长的下界。推导与证明见附录 D。

$L(D)$ 是方案 $D$ 自身的必要耗时下界：$L_{\mathrm{load}}$ 仅取决于输入图与 $N$，$L_{\mathrm{bw}}$、$L_{\mathrm{relay}}$ 随方案变化——方案自身的搬运量与依赖链安排同样抬高其下界。因此 $L(D)$ 不构成对优化问题最优值的估计，也不参与搜索过程；其用途是刻画方案受哪一类机制瓶颈约束（如所得方案贴近带宽下界时，进一步改进须从减少搬运入手），相关分析见 4.6 节。

\subsection{结构感知构造与段单元划分}
\pending{TODO（4.3，待方法侧窗口起草）：链分解（弱连通分量）与拓扑段单元的构造；辫状图自适应段划分（链模式/段模式）。段批合法性引理：批合并当且仅当全局排名相邻且同核，等价于商图无环（附构造证明，成品 tex 在算法侧，按 D018 输入契约挂接）。batch×层窗口重标：由单元指派生成 node\_to\_subgraph 与 core\_schedules 的映射方式。划分合法性的保证方式（与 4.2.3 节约束的对应）。}

\subsection{学习式优先级指派}
\pending{TODO（4.4，待方法侧窗口起草）：每单元优先级参数（logits）；Gumbel-Top-k 采样将单元到核的指派化为一次性可并行排序。领域化解码器：贪心指派代价（M/V 管道负载、跨核共享张量字节项、P3 摊开奖励项 $\lambda$）——4.2.2 节所述启发式代价的落点。离线训练（REINFORCE）：梯度估计、批内代价标准化与 logit 范数正则；方法出处（ICLR23）。与 4.3 的衔接：采样序到单元指派再到重标为完整调度方案。}

\subsection{求解算法}
\pending{TODO（4.5，待方法侧窗口起草）：在线段总体流程——加载离线参数、结构感知构造、贪心＋采样解码、预算内爬山精修、输出方案（单用例硬预算口径）；作为 2.5 节两段式总路线的本章实例化（分层互引，不重复）。}

\subsubsection{算法总体流程}
\pending{TODO：各模块触发顺序与信息流（引用 F07 两段式求解框架示意）。}

\subsubsection{核心策略与模块衔接}
\pending{TODO：精修邻域（单元搬核/换核、B 邻域、序扰动爬山）；接受准则与修复。}

\subsubsection{算法伪代码及停止条件}
\pending{TODO：伪代码对应在线求解流程（算法 1）；停止条件＝墙钟预算。}

\subsubsection{算法复杂度、计算预算与合规性}
\pending{TODO：预算口径（单用例分钟级、非暴力迭代的题面脚注对照）；离线/在线两段分工说明。4.3、4.4 解释策略，4.5 说明如何组织为完整算法，避免重复展开。}

\subsection{实验结果与分析}
\pending{TODO（4.6，数字源＝合规产线 harvest，D038）：主结果曲线、逐用例附录表、求解时间合规统计、消融矩阵（A–E）、bit-exact 对拍证据（FastEval 一致性验证落点，4.2.4 节引用处）、离线校准上界 gap 指针（详表见 7.2/TBL-011）。}

\subsubsection{实验设置与评价方法}
\pending{TODO：测试集、核数、运行环境及版本；官方配置、指标定义与聚合方法；单核基准和多核对照方法；失败、超时、缺失与兜底的处理；三问共用的实验协议（TBL-009）；FastEval 与官方评估器一致性证据（历史全量对拍记录＋harvest 复评 verify 计数）。}

\subsubsection{1～5 核平均加速比及核数扩展表现}
\pending{TODO：明确单核基准点，结合曲线分析性能随核数变化的情况（F19 核数曲线、TBL-004 主结果表、F08 逐例分布）。}

\subsubsection{Makespan 与总额外数据搬运量}
\pending{TODO：整体结果及用例差异，时间与搬运之间的关系（F02 条件图，数字并入 TBL-004）。}

\subsubsection{基线对比与关键策略消融}
\pending{TODO：消融矩阵 A–E（采样数/学习先验/精修开关/解码器项/多种子）；池线历史阶梯作算法族贡献对照分列呈现、逐段标协议（TBL-010/F12/F21）。}

\subsubsection{求解成本、稳定性与典型案例}
\pending{TODO：算法耗时、评估次数、不利案例；重复运行分析（timing\_stats 预算合规）；方案机制下界 $L(D)$ 的瓶颈诊断使用（4.2.4 节定义处引用；甘特图 F03 待确认）。}

\subsubsection{算法结果小结}
\pending{TODO：已证实结果、适用条件与未解决问题；离线校准上界 gap 指针（TBL-011 见 7.2）。}
